% Bagian: intuitionistic-logic
% Bab: introduction
% Subbagian: constructive-reasoning

\documentclass[../../../include/open-logic-chapter]{subfiles}

\begin{document}

\olfileid[id]{int}{int}{cr}

\section{Penalaran Konstruktif}

Berbeda dari perluasan logika klasik dengan operator modal atau kuantor
orde kedua, logika intuisionistik bersifat ``nonklasik'' karena membatasi
logika klasik. Logika klasik bersifat \emph{nonkonstruktif} dalam berbagai
hal. Logika intuisionistik dimaksudkan untuk menangkap jenis penalaran yang
lebih ``konstruktif'', yang menjadi ciri suatu bentuk matematika konstruktif.
Contoh-contoh berikut dapat menggambarkan beberapa motivasi yang mendasarinya.

Andaikan seseorang mengaku telah menentukan suatu bilangan asli~$n$ dengan
sifat bahwa jika $n$ genap, hipotesis Riemann benar, dan jika $n$ ganjil,
hipotesis Riemann salah. Kabar bagus!{} Benar atau tidaknya hipotesis Riemann
merupakan salah satu pertanyaan terbuka besar dalam matematika, dan orang itu
tampaknya telah mereduksi masalah tersebut menjadi masalah perhitungan, yakni
menentukan apakah suatu bilangan tertentu genap atau bukan.

Berapakah nilai ajaib~$n$? Orang itu menguraikannya sebagai berikut: $n$
adalah bilangan asli yang sama dengan $2$ jika hipotesis Riemann benar, dan
$3$ jika tidak.

Dengan marah, Anda menuntut uang Anda kembali. Dari sudut pandang klasik,
uraian di atas memang menentukan satu nilai unik bagi $n$; tetapi yang
sesungguhnya Anda inginkan ialah nilai $n$ yang diberikan secara
\emph{eksplisit}.

Sebagai contoh lain yang mungkin tidak terlalu dibuat-buat, pertimbangkan
pertanyaan berikut. Kita tahu bahwa suatu bilangan irasional dapat dipangkatkan
dengan bilangan rasional dan menghasilkan bilangan rasional. Misalnya,
$\sqrt{2}^2 = 2$. Yang kurang jelas ialah apakah suatu bilangan irasional dapat
dipangkatkan dengan bilangan \emph{irasional} dan menghasilkan bilangan
rasional. Teorema berikut memberikan jawaban afirmatif:

\begin{thm}
Terdapat bilangan irasional $a$ dan $b$ sedemikian sehingga $a^b$ rasional.
\end{thm}

\begin{proof}
Pertimbangkan $\sqrt{2}^{\sqrt{2}}$. Jika bilangan ini rasional, kita selesai:
kita dapat mengambil $a = b = \sqrt{2}$. Jika tidak, bilangan itu irasional.
Selanjutnya,
\[
(\sqrt{2}^{\sqrt{2}})^{\sqrt{2}} = \sqrt{2}^{\sqrt{2} \cdot
  \sqrt{2}} = \sqrt{2}^2 = 2,
\]
yang rasional. Jadi, dalam kasus ini, ambil $a$ sebagai
$\sqrt{2}^{\sqrt{2}}$, dan ambil $b$ sebagai~$\sqrt 2$.
\end{proof}

Apakah ini merupakan bukti yang valid? Kebanyakan matematikawan merasa
demikian. Namun, sekali lagi, ada sesuatu yang agak tidak memuaskan di sini:
kita telah membuktikan keberadaan sepasang bilangan real dengan sifat tertentu
tanpa dapat mengatakan pasangan bilangan \emph{mana} yang dimaksud. Hasil yang
sama dapat dibuktikan dengan cara yang \emph{memberikan} pasangan $a$, $b$ di
dalam bukti itu sendiri: ambil $a = \sqrt{3}$ dan $b = \log_3 4$. Maka
\[
a^b = \sqrt{3}^{\log_3 4} = 3^{1/2 \cdot \log_3 4} = (3^{\log_3
  4})^{1/2} = 4^{1/2}= 2,
\]
karena $3^{\log_3 x} = x$.

Logika intuisionistik dirancang untuk menangkap jenis penalaran yang tidak
memperbolehkan langkah seperti dalam bukti pertama. Membuktikan keberadaan
suatu $x$ yang memenuhi~$!A(x)$ berarti kita harus memberikan suatu~$x$
tertentu dan bukti bahwa nilai itu memenuhi~$!A$, seperti dalam bukti kedua.
Membuktikan bahwa $!A$ atau $!B$ berlaku mengharuskan kita dapat membuktikan
salah satunya.

Secara formal, logika intuisionistik diperoleh dengan membatasi suatu sistem
!!a{derivation} bagi logika klasik dengan cara tertentu. Dari sudut pandang
matematis, semua itu hanyalah sistem deduktif formal, tetapi, sebagaimana telah
dicatat, sistem-sistem tersebut dimaksudkan untuk menangkap suatu jenis
penalaran matematis. Orang dapat menganggapnya sebagai jenis penalaran yang
dibenarkan oleh pandangan filosofis tertentu tentang matematika (seperti
intuisionisme Brouwer); orang dapat menganggapnya sebagai jenis penalaran
matematis yang lebih ``konkret'' dan memuaskan (sejalan dengan konstruktivisme
Bishop); dan orang dapat memperdebatkan apakah uraian formal tersebut
menangkap motivasi informalnya. Namun, apa pun pendirian filosofis kita, kita
dapat mengkaji logika intuisionistik sebagai logika yang disajikan secara
formal; dan apa pun alasannya, banyak ahli logika matematika tertarik untuk
melakukannya.

\end{document}
